%!TEX program = lualatex
\documentclass[11pt]{article}

% ---------- Font setup ----------
\usepackage{fontspec}
\usepackage{unicode-math}
\setmainfont{Libertinus Sans}[Scale=1.1]
\setmathfont{Libertinus Math}[Scale=1.1]

% ---------- Page layout ----------
\usepackage[a4paper,margin=1in]{geometry}

% ---------- Line spacing ----------
\usepackage{setspace}
\setstretch{0.96}
\usepackage[colorlinks=true, urlcolor=blue]{hyperref}

% ---------- Packages ----------
\usepackage{tikz-cd}

% ---------- Custom commands ----------
\newcommand{\Lan}{\operatorname{Lan}}
\newcommand{\Ran}{\operatorname{Ran}}
\newcommand{\Ho}{\operatorname{Ho}}
\newcommand{\id}{\operatorname{id}}
\newcommand{\Top}{\mathbf{Top}}
\newcommand{\ChAb}{\mathbf{Ch}_\bullet(\mathbf{Ab})}
\newcommand{\calC}{\mathcal{C}}
\newcommand{\calD}{\mathcal{D}}

% ---------- Title ----------
\title{
    Reading Group on Categorical Homotopy Theory
    \bigbreak
    \large Introduction and outline for the first 4 meetings
}
\author{}
\date{\today}

\begin{document}

\maketitle

\section{To invert weak equivalences}
In the category $\Top$ of topological spaces,
homotopy equivalence is a good equivalence relation,
so we can define the homotopy category $\Ho\Top$,
whose objects are all the topological spaces
but morphisms are equivalence classes of continuous maps up to homotopy equivalence.
Sometimes working in $\Ho\Top$ is easy since every homotopy equivalence becomes an isomorphism. 

There are also other interesting notions of ``weak equivalences'' that are not actually equivalences.
Again, in the category $\Top$ of topological spaces, we say a map $f : X \to Y$ is a
\href{https://ncatlab.org/nlab/show/weak+homotopy+equivalence}{weak homotopy equivalence} if
$$\pi_n(f) : \pi_n(X) \to \pi_n(Y)$$
is an isomorphism for $n\ge 1$ and a bijection when $n = 0$.
Clearly homotopy equivalence is a weak equivalence, and we have the classical Whitehead's theorem
(which is not provable in HoTT!) saying that,
every weak homotopy equivalence between CW-complexes is a homotopy equivalence.
However, weak homotopy equivalences are generally not invertible
(see \href{https://ncatlab.org/nlab/show/pseudocircle}{here} for a counterexample).

Another important example comes from the category of
\href{https://ncatlab.org/nlab/show/chain+complex}{chain complexes}.
\href{https://ncatlab.org/nlab/show/quasi-isomorphism}{Quasi-isomor-phisms},
as the name might indicate, are also generally not invertible.

What we can do is formally add an invertible arrow to theses weak equivalences.
We consider the following construction of $\Ho\Top_{whe}$.
The objects are all the topological spaces.
The morphisms bewtween $X$ and $Y$ are the equivalence classes of zig-zags
$$X \leftarrow Z_1 \rightarrow \cdots \leftarrow Z_{n} \to Y,$$
suject to
\begin{enumerate}
    \item all the arrows backwards are weak homotopy equivalences.
    \item identity arrows can be removed and then composed.
\end{enumerate}

Then a weak homotopy equivalence $X \stackrel{w}{\to} Y$ has an inverse given by
$$X \stackrel{w}{\leftarrow} X \stackrel{\id}{\to} Y.$$

However, this construction may not give a locally small category.
That is, given a category $\mathcal C$ and a class of morphisms $W$ we want to want to invert,
the resulting $\Ho_W \mathcal C(X,Y)$ may be a proper class
(or in type theory, this means the type would lie in the next universe level).

This is one of the reasons why we need
\href{https://ncatlab.org/nlab/show/model+category}{model category}.
With auxiliary data of fibrations and cofibrations,
we can ensure the homotopy category we obtain is indeed locally small.

Given two categories $\calC$ and $\calD$ and a functor $F : \calC \to \calD$,
we may want to extend $F$ to $\Ho\cal C \to \Ho\cal D$.
That is, we want $\tilde F$ such that the following diagram commutes
\[\begin{tikzcd}
\calC \arrow[r,"F"] \arrow[d,"\gamma"'] & \calD \arrow[d,"\gamma"] \\
\Ho \calC \arrow[r,"\tilde{F}"'] & \Ho \calD
\end{tikzcd}\]

This is easy if $F$ preserves weak equivalences.
But unfortunately, this is most of the time not the case.
For example, the following two diagrams are (component-wise) homotopy equivalent

\[
\begin{tikzcd}
    D ^ 1 \arrow[d] & S ^ 0 \arrow[l] \arrow[r] \arrow[d] & D ^1 \arrow[d] \\
    * & S ^ 0 \arrow[l] \arrow[r] & *
\end{tikzcd}
\]
But the colimit (a functor $\Top^{\bullet \leftarrow \bullet \rightarrow} \to \Top$) of the first row
is $S^1$, while the colimit of the second row is $*$. 

So we have to compromise by considering a weaker notion of extension,
i.e., Kan extension of the functor $F$ along $\gamma: \calC \to \Ho\calC$.
This gives us the so-called derived functors.
In particular, homotopy co/limits are then the derived functors of (strict) co/limits.

\section{``Weak equivalences are all that matter''}
Traditionally, all the constructions mentioned above take place in model categories.
But a generalization is possible: we can introduce a framework
for constructing derived functors between categories equipped with a reasonable notion of
weak equivalence. As Riehl said in the preface:

\begin{quote}
``The weak equivalences are all that matter in abstract homotopy theory, showing
that particular notions of cofibrations/fibrations and cofibrant/fibrant objects are irrelevant
to the construction of derived functors.''
\end{quote}

This generalization allows us to produce derived functors where
appropriate model structures for the diagrams do not necessarily exist.

Also, this philosophy could be inspiring for type theorists.
In type theories, we rarely have a nice model structure,
but we always have a nice notion of equivalence.
So we would at least hope studying this book can help me understand
homotopy co/limits in type theories like HoTT,
or compare homotopy co/limits and strict co/limits in 2LTT.
In 2.2, the book introduces this notion:
a left deformation on $\calC$ is an endofunctor $Q$
with a weak equivalence $Q \Rightarrow 1$.
Doesn't it resemble modality?
We hope they can be somehow connected.

\section{An outline for the first 4 weeks}

We hope to cover most of the first 3 chapters in the first 4 weeks.

\subsection{Week 1: Kan extensions}
\begin{quote}
``All concepts are Kan extensions.''
\end{quote}

Yeah, we will introduce Kan extensions, the last part in Riehl's introductory book on category theory:
\begin{enumerate}
    \item the defintion
    \item a computation in nice settings
    \item nice Kan extensions that we only care
    \item why ``all concepts are Kan extensions.''
    \item time permitted, we might say a little bit about its application in simplicial sets.
\end{enumerate}

\subsection{Week 2: Homtopical category}

The defintion of homotopical category. If people are interested, we may focus on the examples.

\subsection{Week 3: Derived functors}

Similarly, we should focus on examples. It may be nice to introduce a little homological algebras.

\subsection{Week 4: Enriched category}

The idea of enrichement is actually quite accessible.
Again, I think we should focus on examples.


\end{document}